\subsection{THETA3-C24: non-CM Weil-type sources in every even dimension}\label{C24}

The source problem in C13 can be made concrete in every even dimension.
The following construction supplies non-CM Jacobians with a balanced
$\Q(i)$-action, all as quotients of Albaneses in the single arithmetic
ball covering tower of C15--C16. It does not assert that their Weil
classes are independent of divisor classes, or that the ball surface
already supplies a nonzero projection onto their Weil spaces.

\paragraph{Statement.}
Let $n\geq2$, and let $M$ be the least power of two satisfying
$M^2\geq8n+4$. There are a smooth projective curve $H_n$ over
$\overline\Q$ of genus $2n$, a principally polarized abelian variety
$A_n=J(H_n)$, and a connected finite \'etale cover $X_n\to Y$ of
Schoen's fixed arithmetic ball surface, such that
\[
 \deg(X_n/Y)\leq196M^2<6272n+3136,\qquad
 \operatorname{Alb}(X_n)\twoheadrightarrow A_n.
 \tag{C24a}
\]
The variety $A_n$ is not of CM type. It has an action of $K=\Q(i)$
whose two embeddings each have multiplicity $n$ on
$H^0(A_n,\Omega^1)$, and whose Rosati involution is complex
conjugation. Thus
\[
 W_K(A_n)=\bigwedge_K^{2n}H^1(A_n,\Q)
\]
is an actual rational two-dimensional Hodge subspace of type $(n,n)$.
The construction gives a place above seven where the geometric
semistable toric rank of $A_n$ is exactly two. In dimension four
one may take $M=8$, so the covering degree in (C24a) is at most
$12544$.

\paragraph{Even-level torsion transfer.}
The torsion-branch lemma C16 remains valid for every positive integer
$N$, with precisely the same degree bounds. For even $N$ its branch
set has cardinality
\[
 |\pi(E[N])|=(N^2+4)/2.
 \tag{C24b}
\]
Indeed, multiplication by $N$ is \'etale in characteristic zero
regardless of parity. In the cover
$C_N=C\times_{E,[N]}E$, the map $C_N\to E$ is ramified at every
point of $E[N]$, because $C\to E$ is ramified at its origin.
Thus every ramification index of the composite to
$E/\{\pm1\}$ above $\pi(E[N])$ is even, including the indices
four at points fixed by $[-1]$. The local normalization
$u^2=v^{2e}$ used in C16 is unchanged. It gives a connected
finite \'etale cover of $C$ of degree at most $dN^2$ dominating
any degree-$d$ curve cover with ramification at most two and
branch contained in this set. The four fixed points of $[-1]$
on $E[N]$ give (C24b) by counting its orbits. Applying C15 and
Schoen's degree-two cover gives a ball cover of degree at most
$2dN^2$; for a hyperelliptic cover this is $4N^2$.

\paragraph{The torsion points and the equation.}
Use the elliptic curve from C16,
\[
 E:\quad y^2=x^3-30x-25,
 \qquad T=(-5,0),\qquad c=3\sqrt5.
\]
Fix embeddings of $\overline\Q$ in $\C$ and in
$\overline\Q_7$. The curve has good reduction at seven:
its displayed discriminant is $1458000$, a $7$-adic unit.
Choose a nonzero point $R\in E[7](\overline\Q)$ reducing
to the identity at this place. Such a point exists. There are
$49$ geometric points of $E[7]$ in characteristic zero,
whereas the geometric seven-torsion in the special elliptic
curve has at most seven points: the differential of $[7]$
vanishes, so its inseparable degree is at least seven and
its separable degree is at most seven. Properness gives
reduction of all these torsion points, and its kernel therefore
contains a nonzero point. In particular $R$ belongs to the
formal group at the chosen place.

On the Kummer line set
\[
 z(P)=\frac{x(P)+5-c}{x(P)+5+c}.
 \tag{C24c}
\]
The addition law gives
\[
 x(P+T)+5=\frac{45}{x(P)+5},\qquad
 z(P+T)=-z(P),\qquad z(-P)=z(P).
 \tag{C24d}
\]
For example the first identity follows by using slope
$y(P)/(x(P)+5)$ in the chord from $P$ to $T$ and using
$x^3-30x-25=(x+5)(x^2-5x-5)$. The two fixed points
of the involution on the Kummer line have coordinates $z=0,\infty$;
their preimages satisfy $2P=T$ and belong to $E[4]$.

The four-element group generated by $P\mapsto-P$ and
$P\mapsto P+T$ acts on $E[M]$. Exactly eight points have
nontrivial stabilizer: the four in $E[2]$ and the four satisfying
$2P=T$. These two sets are disjoint, and translation by $T$
has no fixed point. Every remaining orbit has four elements.
There are therefore
\[
 (M^2-8)/4\geq2n-1
\]
such orbits. Choose representatives $U_1,\ldots,U_{2n-1}$
of distinct ones. Put
\[
 P_1=U_1,\qquad P_2=U_1+R,\qquad
 P_j=U_{j-1}\quad(3\leq j\leq2n),\qquad a_j=z(P_j).
 \tag{C24e}
\]
These points lie in $E[7M]$. They give pairwise distinct
branch values
\[
 0,\infty, a_1,-a_1,\ldots,a_{2n},-a_{2n}.
 \tag{C24f}
\]
For the $U_j$ this is precisely the orbit condition.
The point $U_1+R$ cannot belong to the four-element orbit
of any $U_j$: such an equality would put the nonzero
seven-torsion point $R$ in $E[M]$, contrary to $(7,M)=1$.
It likewise cannot be in either excluded set. The branch
values in (C24f) all belong to $z(E[7M])$.

Define the smooth projective model of
\[
 H_n:\qquad v^2=z\prod_{j=1}^{2n}(z^2-a_j^2).
 \tag{C24g}
\]
The distinctness just proved gives a smooth hyperelliptic curve
with $4n+2$ branch points, hence genus $2n$.
Even-level transfer with $N=7M$ gives
$\deg(X_n/Y)\leq4(7M)^2=196M^2$ and the asserted
Albanese quotient. Minimality of $M$ gives
$(M/2)^2<8n+4$, proving the strict final bound in (C24a).

\paragraph{The balanced action.}
The automorphism
\[
 \alpha:(z,v)\longmapsto(-z,iv)
\]
has square equal to the hyperelliptic involution. On the
Jacobian the latter acts as $[-1]$, because for a Weierstrass
base point the divisors $P+\iota(P)$ are all linearly equivalent
to twice that base point. Thus $i\mapsto\alpha^*$ gives
a faithful $\Q(i)$-action on $A_n$. Its pullback on the
holomorphic basis is
\[
 \alpha^*\left(z^j\frac{dz}{v}\right)
   =i(-1)^j z^j\frac{dz}{v},\qquad 0\leq j<2n.
\]
There are $n$ even and $n$ odd exponents, giving the claimed
signature. A curve automorphism preserves the canonical
principal polarization of its Jacobian; hence its Rosati
adjoint is its inverse. Here $\alpha^{-1}=-\alpha$ on
$H^1$, so this involution restricts to complex conjugation
on $K$. The two $K$-eigenspaces of $H^1_\C$ each have
Hodge numbers $(n,n)$. Their top exterior powers therefore
have type $(n,n)$, as asserted for $W_K(A_n)$.

\paragraph{A two-node reduction certifying that the source is not CM.}
Pass to a finite extension of $\Q_7$ containing all the
displayed points and constants, with valuation ring $\mathcal O$.
Prime-to-seven torsion reduces injectively on the good model
of $E$. Moreover $c$ and $2c$ are units. Consequently
$a_1,a_3,\ldots,a_{2n}$ are units with pairwise distinct
reductions up to sign, none equal to zero or infinity.
The same assertions hold for $a_2$ except that
\[
 \bar a_2=\bar a_1=:a\ne0,\qquad a_2\ne\pm a_1.
\]
Write $b_j=\bar a_j$ for $j\geq3$. A proper flat double
cover of $\mathbf P^1_{\mathcal O}$ is obtained by gluing
(C24g) to its chart $u=1/z$, $w=v/z^{2n+1}$, whose equation is
\[
 w^2=u\prod_{j=1}^{2n}(1-a_j^2u^2).
\]
Each chart is free of rank two over the indicated affine
line. They give a finite cover of $\mathbf P^1$, so it is
proper and flat. The special fibre is
\[
 v^2=z(z^2-a^2)^2\prod_{j=3}^{2n}(z^2-b_j^2).
 \tag{C24h}
\]
It is geometrically irreducible: the squarefree part on the
right has the simple root $z=0$ and is not a square.
Its only singularities are the two points $z=a,-a$, each
an ordinary node. At each, the remaining factor is a unit
and $2$ is invertible, so the tangent cone is the product
of two distinct linear factors. The points at zero and
infinity and the other branch points are smooth.
The normalization is
\[
 v_0^2=z\prod_{j=3}^{2n}(z^2-b_j^2),
 \qquad v_0=v/(z^2-a^2),
\]
of genus $2n-2$. In particular the fibre is a stable
irreducible curve with two nodes.

For completeness, allowing a singular total space does not
alter the toric-rank calculation. In the completed local
ring near either node, extract a square root of the unit
factor and put $\delta=a_2-a_1$. The remaining equation
is equivalent to
\[
 XY=\delta^2/4
\]
up to a unit, by centering $z$ at $(a_1+a_2)/2$ and
factoring the difference of two squares. A finite unramified
extension suffices to split the two branches. If
$v(\delta)=e$, resolution of $XY=\pi^{2e}\cdot\mathrm{unit}$
inserts a chain between those two branches. This can also
be seen inductively in the two blow-up charts, which reduce
the exponent by one. In the dual graph it only subdivides
the corresponding loop. After any finite ramified base
extension the exponent is multiplied by its ramification
index and the same description applies. The geometric
dual graph therefore has first Betti number two over
every such extension.

Apply the N\'eron--Picard identification in SGA~7~I,
Expos\'e~IX, Theorem~12.1(d), formulas (12.1.10)--(12.1.12), printed pp.~467--468,
to a regular semistable resolution. The torus of the
special-fibre Picard variety has character lattice the
first homology of this dual graph, as in IX~12.3,
especially (12.3.6)--(12.3.7), printed pp.~471--472. Equivalently, gluing line
bundles across the two independent nodes gives two
independent $\mathbf G_m$ parameters. Hence the Jacobian
has toric rank exactly two, and cannot acquire good
reduction after any finite extension. Serre--Tate,
\emph{Good reduction of abelian varieties} (1968),
Theorem~6(a), printed p.~503, gives potential good
reduction for CM abelian varieties. For a product of
simple CM factors the same conclusion follows after a
common finite extension and isogeny, so it applies to
the general meaning of CM type used here. Thus $A_n$
is not CM. \hfill$\square$

\paragraph{What this advances, and what remains.}
For every $n\geq2$ there is now an explicit finite torsion
prescription for a non-CM source with the precise balanced
quadratic-field action required by C13, within one fixed
arithmetic ball tower. All of its weight-one realization
comes from the actual Albanese quotient (C24a). The
projector in C13 consequently applies to this specific
$W_K(A_n)$, without an additional source-existence
hypothesis. Geometric simplicity, an exact endomorphism
ring, and a non-divisor Weil class have not been inferred
from the nodal calculation. The remaining cycle problem
is still to construct an $n$-dimensional correspondence
with nonzero projection to this Weil space. Taking
products of the supplied curve map produces the
Lefschetz classes described in C14, and does not resolve
that projection problem when the Weil space is
exceptional. For $n=2$, the cited Markman theorem already
gives algebraicity; the new conclusion here is the
specified non-CM source in the ball tower.

\framework{The Albanese maps and rational Weil projector are actual
correspondences in the setting of Grothendieck's cote~44,
scans~7--8, and cote~16, scans~43--46. The new non-CM
certificate takes place directly in his theory of N\'eron
models and monodromy: SGA~7~I, IX~12.1 and IX~12.3,
with the dual-graph torus and its persistent rank. It
supplies non-CM weight-one motives and their specified
Weil submotives modulo homological equivalence; it does
not identify an unknown Hodge class with an effective
special-fibre intersection.}
